%% =====================================================================
%%  Unit 1 — Hello!  (Saluer et se présenter)
%%  Pistes : 1.01–1.04 dialogues · 1.05–1.07 vocabulaire · 1.08 grammaire
%%           1.09 prononciation · 1.10–1.13 exercices (textes en annexe)
%% =====================================================================

% Illustration du bandeau d'ouverture
\newcommand{\otIllusUnI}{%
  \foreach \mot/\x/\y/\t/\o in {Hello!/104/86/15/0.10,Hi!/190/90/22/0.09,Hey!/122/14/13/0.08,
      Good morning/150/8/11/0.08,Welcome!/196/22/12/0.07,Bye!/100/50/12/0.07}{%
    \node[text=white,opacity=\o,font=\otBold\fontsize{\t}{\t}\selectfont,anchor=center] at (\x mm,\y mm) {\en{\mot}};}
  \fill[white,opacity=0.07] (160mm,44mm) circle (40mm);
  \otDessinAvatar{jack}{(137mm,40mm)}{15mm}{otNavyDark!70!otNavy}
  \otDessinAvatar{lea}{(182mm,47mm)}{17.5mm}{otGreen!35!otNavy}
  \otBulle{(128mm,80mm)}{south}{Hi! I'm Jack.}{white}{otNavy}{(134mm,58mm)}
  \otBulle{(181mm,84mm)}{south}{Hello! I'm Léa.}{otGreen}{white}{(181mm,67mm)}}

\unite{1}{Hello!}{Saluer et se présenter}{otIllusUnI}

\situation{Épisode 1}{Lundi, 9 heures du matin, à New York. Léa Martin, graphiste de 27 ans
venue de Lyon, commence aujourd'hui son nouveau travail chez \en{\textit{Bright Ideas}}, un
studio de design de Manhattan. Premier jour, premières rencontres… et premiers mots d'anglais !}

\begin{objectifs}
\item dire bonjour et au revoir, à chaque moment de la journée ;
\item vous présenter et demander son nom à quelqu'un ;
\item demander comment ça va, et répondre ;
\item utiliser les pronoms sujets et le verbe \en{\textbf{be}} (être) ;
\item prononcer le \en{\textbf{h}} anglais, qui ne s'entend jamais en français.
\end{objectifs}

\auprogramme
  {Saluer, prendre congé, se présenter, prendre des nouvelles}
  {Les pronoms sujets\par\vspace{0.6mm}Le verbe \en{be} au présent}
  {Salutations, politesse, présentations : 36 expressions}
  {Le son \phon{h} : un souffle}
  {\en{Hi}, \en{hello} ou \en{good morning} ?\par\vspace{0.6mm}\en{Mr.}, \en{Mrs.}, \en{Ms.}}

\semaine{casque/otNavy/Dialogues/Écouter et lire les deux scènes,
  bulle/otGreenDark/Vocabulaire/Apprendre les 36 expressions,
  blocs/otViolet/Grammaire/Pronoms et verbe be,
  sons/otOrangeDark/Sons et culture/Le h anglais et les usages,
  crayon/otNavy/Exercices/12 exercices corrigés,
  coche/otGreenDark/Bilan/L'essentiel et le test express}

\clearpage
%% ---------------------------------------------------------------------
\rubrique{ecouter}{Premier jour à New York}

\begin{astuce}[Comment travailler un dialogue]
\begin{etapes}
\item Écoutez la \textbf{version lente} sans lire : essayez de reconnaître des mots.
\item Réécoutez en suivant le texte anglais, puis lisez la traduction.
\item Écoutez la \textbf{version normale} et répétez chaque réplique à voix haute, en imitant la voix.
\end{etapes}
\end{astuce}

\begin{dialogue}{1}{Hi, I'm Jack!}{9 heures, à l'accueil du studio \en{\textit{Bright Ideas}}.
  Léa fait la connaissance de deux collègues : Jack, américain, et Emma, britannique.}{1.01}{1.02}
\dit{jack}{\hl{Hi!} \hl{I'm} Jack.\note{1}}{Salut ! Moi, c'est Jack.}
\dit{lea}{\hl{Hello}, Jack. I'm Léa. Léa Martin.}{Bonjour, Jack. Moi, c'est Léa. Léa Martin.}
\dit{jack}{\hl{Nice to meet you}, Léa.\note{2}}{Enchanté, Léa.}
\dit{lea}{Nice to meet you \hl{too}.}{Enchantée moi aussi.}
\dit{jack}{\hl{Welcome} to Bright Ideas!}{Bienvenue chez \textit{Bright Ideas} !}
\dit{lea}{\hl{Thank you}!}{Merci !}
\dit{jack}{And \hl{this is} Emma.\note{3}}{Et voici Emma.}
\dit{emma}{Hello, Léa! Nice to meet you.}{Bonjour, Léa ! Ravie de te rencontrer.}
\dit{lea}{Hi, Emma. Nice to meet you too.}{Salut, Emma. Enchantée moi aussi.}
\dit{emma}{\hl{How are you?}\note{4}}{Comment ça va ?}
\dit{lea}{\hl{I'm fine, thanks.} \hl{And you?}}{Ça va bien, merci. Et toi ?}
\dit{emma}{\hl{I'm very well}, thank you.}{Très bien, merci.}
\end{dialogue}

\begin{comprendre}
\pointnote{1} \en{\textbf{I'm Jack.}} \en{I'm} est la forme courte de \en{I am} (« je suis »).
  Pour se présenter, on dit simplement \en{I'm} + prénom : c'est notre « Moi, c'est… » ou
  « Je m'appelle… ».
\pointnote{2} \en{\textbf{Nice to meet you.}} Mot à mot : « agréable de te rencontrer ». C'est
  la formule de la première rencontre, l'équivalent de « Enchanté(e) ». On répond
  \en{Nice to meet you too.} (\en{too} = aussi).
\pointnote{3} \en{\textbf{This is Emma.}} \en{This is…} (« ceci est… ») sert à présenter
  quelqu'un : « Voici Emma », « Je te présente Emma ».
\pointnote{4} \en{\textbf{How are you?}} « Comment ça va ? » On répond brièvement :
  \en{I'm fine, thanks. And you?} Remarquez que \en{you} veut dire à la fois « tu » et
  « vous » : en anglais, pas besoin de choisir !
\end{comprendre}

\begin{dialogue}{2}{Good morning, Mr.\ Carter}{9 h 30. Ben Carter, le directeur du studio,
  reçoit Léa dans son bureau. Puis, à 18 heures, c'est l'heure de rentrer.}{1.03}{1.04}
\dit{carter}{\hl{Good morning}. \hl{My name is} Ben Carter.}{Bonjour. Je m'appelle Ben Carter.}
\dit{lea}{Good morning, \hl{Mr.\ Carter}. I'm Léa Martin.\note{5}}{Bonjour, monsieur Carter. Je suis Léa Martin.}
\dit{carter}{Ah, Léa! Welcome to New York. \hl{Please}, call me Ben.\note{6}}{Ah, Léa ! Bienvenue à New York. Je vous en prie, appelez-moi Ben.}
\dit{lea}{\hl{Thank you very much}\ldots{} Ben.}{Merci beaucoup… Ben.}
\dit{carter}{How are you today?}{Comment allez-vous aujourd'hui ?}
\dit{lea}{I'm very well, thank you. And you?}{Très bien, merci. Et vous ?}
\dit{carter}{\hl{I'm good}, thanks!}{Bien, merci !}
\didascalie{18 heures. Léa quitte le studio.}
\dit{jack}{\hl{Bye}, Léa! \hl{See you tomorrow!}\note{7}}{Salut, Léa ! À demain !}
\dit{lea}{\hl{Goodbye}, Jack. \hl{Have a nice evening!}}{Au revoir, Jack. Bonne soirée !}
\dit{jack}{Thanks, \hl{you too}!}{Merci, toi aussi !}
\end{dialogue}

\begin{comprendre}
\pointnote{5} \en{\textbf{Mr.\ Carter}} \en{Mr.} (prononcé \phon{ˈmɪstɚ}) veut dire « monsieur » et
  se place toujours devant le \textbf{nom de famille} : \en{Mr.\ Carter}, jamais \en{Mr.\ Ben}.
\pointnote{6} \en{\textbf{Please, call me Ben.}} « Je vous en prie, appelez-moi Ben. » Aux
  États-Unis, on passe très vite au prénom, même avec son directeur. \en{Please} = « s'il te
  plaît, s'il vous plaît ».
\pointnote{7} \en{\textbf{See you tomorrow!}} Mot à mot : « je te vois demain » = « À demain ! ».
  Pour souhaiter une bonne soirée : \en{Have a nice evening!} On répond \en{You too!}
  (« Toi aussi ! / Vous aussi ! »).
\end{comprendre}

%% ---------------------------------------------------------------------
\rubrique{mots}{Les mots de l'unité}

Voici les 36 expressions de l'unité, classées par situation. La colonne
« Prononciation » utilise l'alphabet phonétique, présenté page~\pageref{lim-sons} :
\phon{ˈ} marque la syllabe forte et \phon{ː} une voyelle longue. Mais la vraie référence,
c'est l'audio : écoutez, répétez, réécoutez.

\begin{vocabulaire}{Dire bonjour et au revoir}{1.05}
\categorie{Saluer}
\mot{hello}{həˈloʊ}{bonjour (à toute heure)}
\mot{hi}{haɪ}{salut, bonjour}
\mot{hey}{heɪ}{salut \registre{familier}}
\mot{good morning}{ɡʊd ˈmɔːrnɪŋ}{bonjour (le matin)}
\mot{good afternoon}{ɡʊd ˌæftɚˈnuːn}{bonjour (l'après-midi)}
\mot{good evening}{ɡʊd ˈiːvnɪŋ}{bonsoir (en arrivant)}
\categorie{Prendre congé}
\mot{goodbye}{ɡʊdˈbaɪ}{au revoir}
\mot{bye}{baɪ}{salut, au revoir}
\mot{see you later}{ˌsiː juː ˈleɪtɚ}{à plus tard}
\mot{see you tomorrow}{ˌsiː juː təˈmɑːroʊ}{à demain}
\mot{have a nice day}{ˌhæv ə ˌnaɪs ˈdeɪ}{bonne journée}
\mot{have a nice evening}{ˌhæv ə ˌnaɪs ˈiːvnɪŋ}{bonne soirée}
\mot{good night}{ɡʊd ˈnaɪt}{bonne nuit, bonne soirée (en partant tard, ou avant de dormir)}
\end{vocabulaire}

\otIntertitre{Repère visuel · Quelle salutation, à quelle heure ?}
\noindent\begin{tikzpicture}[x=6.75mm]
  % bande de la journée
  \fill[otNavyDark!80] (0,0) rectangle (5,7mm);
  \fill[otOrange!55] (5,0) rectangle (12,7mm);
  \fill[otOrange!85] (12,0) rectangle (18,7mm);
  \fill[otViolet!70] (18,0) rectangle (24,7mm);
  \foreach \h in {0,6,12,18,24}{\draw[otMuted,line width=0.5pt] (\h,0) -- (\h,-1.4mm);
    \node[below,inner sep=1.4mm,font=\otRegular\fontsize{7.4}{9}\selectfont,text=otMuted] at (\h,-0.6mm) {\h~h};}
  \node[font=\otMedium\fontsize{7.6}{9}\selectfont,text=white] at (2.5,3.5mm) {la nuit};
  \node[font=\otSemi\fontsize{9.4}{11}\selectfont,text=otInk] at (8.5,3.5mm) {\enb{Good morning}};
  \node[font=\otSemi\fontsize{9.4}{11}\selectfont,text=otInk] at (15,3.5mm) {\enb{Good afternoon}};
  \node[font=\otSemi\fontsize{9.4}{11}\selectfont,text=white] at (21,3.5mm) {\enb{Good evening}};
  \otDessinIcone{soleil}{otOrangeDark}{white}{6mm}{(8.5,12mm)}
  \otDessinIcone{soleil}{otOrange}{white}{6mm}{(15,12mm)}
  \otDessinIcone{lune}{otViolet}{white}{5.2mm}{(21,12mm)}
  % à toute heure
  \draw[otGreenDark,line width=0.8pt,{Stealth[length=2mm]}-{Stealth[length=2mm]}] (0,-11mm) -- (24,-11mm);
  \node[fill=white,inner xsep=2mm,font=\otSemi\fontsize{9}{11}\selectfont,text=otGreenDark] at (12,-11mm)
    {\enb{Hello! · Hi!}\enspace{\otRegular\color{otMuted}à toute heure}};
  % good night
  \draw[otViolet,line width=0.6pt,-{Stealth[length=1.8mm]}] (22.6,18.5mm) -- (22.6,8mm);
  \node[anchor=south east,inner sep=0pt,align=right,font=\otRegular\fontsize{8}{10}\selectfont,text=otMuted] at (24,20mm)
    {{\otSemi\color{otViolet}\enb{Good night}} : on le dit en \textbf{partant},\\ jamais en arrivant};
\end{tikzpicture}

\begin{vocabulaire}{Se présenter et prendre des nouvelles}{1.06}
\categorie{Se présenter}
\mot{I'm Léa.}{aɪm ˈleɪə}{Moi, c'est Léa. / Je suis Léa.}
\mot{My name is Jack.}{maɪ ˈneɪm ɪz ˈdʒæk}{Je m'appelle Jack.}
\mot{What's your name?}{ˌwʌts jʊr ˈneɪm}{Comment tu t'appelles ? / Comment vous appelez-vous ?}
\mot{Nice to meet you.}{ˌnaɪs tə ˈmiːt juː}{Enchanté(e).}
\mot{Nice to meet you too.}{ˌnaɪs tə ˌmiːt juː ˈtuː}{Enchanté(e) moi aussi.}
\mot{This is Emma.}{ˈðɪs ɪz ˈemə}{Voici Emma. / Je te présente Emma.}
\mot{Welcome!}{ˈwelkəm}{Bienvenue !}
\categorie{Prendre des nouvelles}
\mot{How are you?}{ˌhaʊ ˈɑːr juː}{Comment ça va ? / Comment allez-vous ?}
\mot{I'm fine, thanks.}{aɪm ˈfaɪn θæŋks}{Ça va bien, merci.}
\mot{I'm good, thanks.}{aɪm ˈɡʊd θæŋks}{Ça va bien, merci. \registre{très courant aux États-Unis}}
\mot{I'm very well, thank you.}{aɪm ˌveri ˈwel ˈθæŋk juː}{Je vais très bien, merci.}
\mot{I'm OK.}{aɪm ˌoʊˈkeɪ}{Ça va.}
\mot{Not bad.}{ˌnɑːt ˈbæd}{Pas mal.}
\mot{And you?}{ənd ˈjuː}{Et toi ? / Et vous ?}
\end{vocabulaire}

\begin{vocabulaire}{La politesse et les titres}{1.07}
\categorie{Être poli}
\mot{please}{pliːz}{s'il te plaît, s'il vous plaît}
\mot{thank you}{ˈθæŋk juː}{merci}
\mot{thanks}{θæŋks}{merci \registre{courant}}
\mot{Thank you very much.}{ˈθæŋk juː ˌveri ˈmʌtʃ}{Merci beaucoup.}
\mot{You're welcome.}{jʊr ˈwelkəm}{De rien. / Je vous en prie.}
\mot{You too!}{ˌjuː ˈtuː}{Toi aussi ! / Vous aussi !}
\categorie{Les titres}
\mot{Mr.\ Carter}{ˈmɪstɚ ˈkɑːrtɚ}{monsieur Carter}
\mot{Mrs.\ Lopez}{ˈmɪsɪz ˈloʊpez}{madame Lopez (femme mariée)}
\mot[Mizz Clarke.]{Ms.\ Clarke}{mɪz ˈklɑːrk}{madame Clarke (sans préciser)}
\end{vocabulaire}

\begin{astuce}[Apprendre ses mots]
Cachez la colonne « Anglais » avec une feuille et retrouvez chaque expression à partir du
français, \textbf{à voix haute}. Recommencez demain, puis dans trois jours : c'est en espaçant
les révisions qu'on retient pour longtemps.
\end{astuce}

%% ---------------------------------------------------------------------
\rubrique{grammaire}{Les pronoms sujets et le verbe be}

\otIntertitre[otViolet]{1 · Les pronoms personnels sujets}
Comme en français, chaque verbe anglais a un sujet. Voici les pronoms qui le remplacent :

\begin{center}\otTableau
\begin{tabular}{@{}>{\enOn\bfseries}l<{\enOff}@{\hspace{2mm}}>{\fontsize{10}{12}\selectfont}l@{\hspace{5mm}}l@{\hspace{14mm}}>{\enOn\bfseries}l<{\enOff}@{\hspace{2mm}}>{\fontsize{10}{12}\selectfont}l@{\hspace{5mm}}l@{}}
\multicolumn{3}{@{}l}{\otTH{Singulier}} & \multicolumn{3}{l@{}}{\otTH{Pluriel}}\\\midrule
I   & \ipa{aɪ}  & je & we & \ipa{wiː} & nous\\
you & \ipa{juː} & tu, vous (politesse) & you & \ipa{juː} & vous (à plusieurs)\\
he  & \ipa{hiː} & il (un homme, un garçon) & they & \ipa{ðeɪ} & ils, elles\\
she & \ipa{ʃiː} & elle (une femme, une fille) & & & \\
it  & \ipa{ɪt}  & il, elle (une chose, un animal) & & & \\\bottomrule
\end{tabular}\end{center}

\begin{regle}[Quatre différences avec le français]
\begin{etapes}
\item \en{\textbf{I}} (je) s'écrit \textbf{toujours} avec une majuscule, même au milieu d'une
  phrase : \en{Jack and I}.
\item \en{\textbf{you}} veut dire « tu », « vous » de politesse et « vous » pluriel : un seul mot
  pour les trois ! La politesse passe par le ton et par des mots comme \en{please},
  \en{thank you} ou \en{Mr.}
\item \en{\textbf{he}} et \en{\textbf{she}} désignent des personnes. Pour une chose, une ville
  ou un animal, on emploie \en{\textbf{it}} : \en{The studio? It's in Manhattan.}
\item \en{\textbf{they}} veut dire « ils » \textbf{et} « elles » : pas de masculin ni de
  féminin au pluriel.
\end{etapes}
\end{regle}

\otIntertitre[otViolet]{2 · Le verbe be (être) au présent}
\en{Be} est le verbe le plus utilisé de l'anglais. Au présent, il n'a que trois formes :
\en{\textbf{am}}, \en{\textbf{is}} et \en{\textbf{are}}. À l'oral, on emploie presque toujours
la \textbf{forme contractée}.

\begin{center}\otTableau
\begin{tabular}{@{}>{\enOn}l<{\enOff}@{\hspace{9mm}}>{\enOn\bfseries}l<{\enOff}@{\hspace{9mm}}l@{\hspace{9mm}}l@{}}
\otTH{Forme pleine} & \otTH{Forme contractée} & \otTH{Prononciation} & \otTH{Français}\\\midrule
I am       & I'm      & \ipa{aɪm}  & je suis\\
you are    & you're   & \ipa{jʊr}  & tu es, vous êtes\\
he is      & he's     & \ipa{hiːz} & il est\\
she is     & she's    & \ipa{ʃiːz} & elle est\\
it is      & it's     & \ipa{ɪts}  & il / elle est, c'est\\\midrule
we are     & we're    & \ipa{wɪr}  & nous sommes\\
you are    & you're   & \ipa{jʊr}  & vous êtes\\
they are   & they're  & \ipa{ðer}  & ils / elles sont\\\bottomrule
\end{tabular}\end{center}

\blocaudio{1.08}{Écoutez et répétez}{La voix lit chaque ligne du tableau : d'abord la forme
pleine, puis la forme contractée. Répétez pendant les pauses.}

\begin{paire}
\begin{regle}[La contraction]
L'apostrophe remplace la lettre qui disparaît :\par\vspace{1mm}
\en{I \textbf{a}m} → \en{I\textbf{'}m}\quad \en{he \textbf{i}s} → \en{he\textbf{'}s}\quad
\en{you \textbf{a}re} → \en{you\textbf{'}re}\par\vspace{1mm}
Forme contractée : à l'oral et dans les messages. Forme pleine : à l'écrit soigné, ou pour insister.
\end{regle}
\begin{regle}[Les emplois de \en{be}]
\exen{I'm Léa.}{l'identité}
\exen{She's very well.}{l'état : « elle va très bien »}
\exen{They're in New York.}{le lieu}
\exen{What's your name?}{mot à mot : « quel est ton nom ? »}
\end{regle}
\end{paire}

\begin{attention}[Les pièges des francophones]
\fauxjuste{I call me Léa.}{My name is Léa. / I'm Léa.}
\fauxjuste{How do you go?}{How are you?}
\fauxjuste{She are Emma.}{She's Emma.}
\fauxjuste{What's you're name?}{What's your name?}
\par\vspace{1mm}{\fontsize{10}{12.6}\selectfont \en{you're} (« tu es ») et \en{your} (« ton, ta,
votre ») se prononcent de la même façon, \phon{jʊr} : seule l'orthographe les distingue.\par}
\end{attention}

\begin{astuce}[Pour ne plus hésiter entre \en{am}, \en{is} et \en{are}]
\en{\textbf{I}} va avec \en{\textbf{am}}. \en{\textbf{He}}, \en{\textbf{she}}, \en{\textbf{it}} — et
tout nom au singulier, comme \en{Jack} ou \en{my name} — vont avec \en{\textbf{is}}.
Tous les autres (\en{you}, \en{we}, \en{they}) vont avec \en{\textbf{are}}.
\end{astuce}

%% ---------------------------------------------------------------------
\rubrique[Le son h]{sons}{Le son \phon{h} : un souffle}

En français, le \textbf{h} ne se prononce jamais : on ne l'entend ni dans « hôtel », ni dans
« homme ». En anglais, il se prononce presque toujours. C'est un \textbf{souffle léger}, comme
celui qui fait de la buée quand on souffle sur ses lunettes : \en{\textbf{h}ello},
\en{\textbf{h}i}, \en{\textbf{h}ow}, \en{\textbf{h}e}.

\begin{tcolorbox}[otBase,colback=white,colframe=otRule,boxrule=0.5pt]
\otTitreBoite{cible}{otOrangeDark}{Le test de la main}{}
\begin{minipage}[c]{0.47\linewidth}\centering
\begin{tikzpicture}
  \otDessinAvatar{lea}{(0,0)}{8.5mm}{otGreenTint}
  \foreach \y in {-2.2,0,2.2}{\draw[otOrange,line width=1pt,line cap=round]
    (10mm,\y mm-2.5mm) .. controls (14mm,\y mm-0.5mm) and (16mm,\y mm-4.5mm) .. (20mm,\y mm-2.5mm)
    .. controls (23mm,\y mm-0.8mm) and (24mm,\y mm-3.5mm) .. (27mm,\y mm-2.5mm);}
  \otDessinMain{(35mm,-2mm)}{10mm}
  \node[font=\otSemi\fontsize{12}{14}\selectfont,text=otOrangeDark] at (17mm,10mm) {\enb{hi} \ipa{haɪ}};
  \node[text width=48mm,align=center,font=\fontsize{9.4}{11.4}\selectfont] at (17mm,-16mm)
    {Vous sentez un petit souffle chaud sur la main.};
\end{tikzpicture}
\end{minipage}\hfill{\color{otRule}\vrule width 0.5pt}\hfill
\begin{minipage}[c]{0.47\linewidth}\centering
\begin{tikzpicture}
  \otDessinAvatar{lea}{(0,0)}{8.5mm}{otGreenTint}
  \otDessinMain{(35mm,-2mm)}{10mm}
  \node[font=\otSemi\fontsize{12}{14}\selectfont,text=otNavy] at (17mm,10mm) {\enb{I} \ipa{aɪ}};
  \node[text width=48mm,align=center,font=\fontsize{9.4}{11.4}\selectfont] at (17mm,-16mm)
    {Vous ne sentez presque rien : pas de \textbf{h}, pas de souffle.};
\end{tikzpicture}
\end{minipage}
\par\vspace{2mm}{\fontsize{10}{12.6}\selectfont Placez votre main à 5\,cm de votre bouche, puis
dites \en{\textit{hi}}, puis \en{\textit{I}}. C'est toute la différence !\par}
\end{tcolorbox}

\begin{paires}{1.09}{Dans chaque paire, un seul son change : le souffle du \textbf{h}.
Écoutez, puis répétez les deux mots en exagérant le souffle.}
\sonpaire{I}{aɪ}{je}{\otH{h}i}{haɪ}{salut}
\sonpaire{and}{ænd}{et}{\otH{h}and}{hænd}{main}
\sonpaire{air}{er}{air}{\otH{h}air}{her}{cheveux}
\sonpaire{eat}{iːt}{manger}{\otH{h}eat}{hiːt}{chaleur}
\sonpaire{old}{oʊld}{vieux}{\otH{h}old}{hoʊld}{tenir}
\sonpaire{art}{ɑːrt}{art}{\otH{h}eart}{hɑːrt}{cœur}
\sonpaire{ear}{ɪr}{oreille}{\otH{h}ere}{hɪr}{ici}
\sonpaire{is}{ɪz}{est}{\otH{h}is}{hɪz}{son, sa (à lui)}
\end{paires}

\begin{attention}[Trois pièges]
\begin{etapes}
\item \textbf{Oublier le h} : on n'entend plus la différence entre \en{hi} et \en{I}, ou entre
  \en{heart} et \en{art}. Dites bien \en{\textbf{h}ello}, \en{\textbf{h}ow are you?}
\item \textbf{En ajouter un} là où il n'y en a pas : à force de faire attention, beaucoup de
  francophones disent « \en{hI'm} » ou « \en{hand you} ». \en{I'm} et \en{and} commencent par
  une voyelle : pas de souffle !
\item \textbf{Les h muets} : quelques mots se prononcent sans h, comme en français :
  \en{hour} \ipa{aʊr} (heure), \en{honest} \ipa{ˈɑːnɪst} (honnête). Ils sont rares :
  vous les apprendrez au fil des unités.
\end{etapes}
\end{attention}

%% ---------------------------------------------------------------------
\rubrique{culture}{Dire bonjour à l'américaine}

\begin{culture}[Du plus familier au plus formel]
\begin{center}
\begin{tikzpicture}
  \draw[otTeal!45,line width=1.6pt,-{Stealth[length=3mm]}] (0,0) -- (150mm,0);
  \foreach \x/\m/\c/\t in {12/Hey!/otTeal!18/otTeal!80!black,48/Hi!/otTeal!40/otTeal!70!black,88/Hello!/otTeal!75/white,130/Good morning./otTeal/white}{%
    \node[fill=\c,text=\t,rounded corners=1.8mm,inner xsep=3mm,inner ysep=1.8mm,
      font=\otSemi\fontsize{10.4}{12}\selectfont] at (\x mm,0) {\enb{\m}};}
  \node[anchor=north,font=\otRegular\fontsize{8}{10}\selectfont,text=otMuted] at (12mm,-5mm) {entre amis};
  \node[anchor=north,font=\otRegular\fontsize{8}{10}\selectfont,text=otMuted] at (48mm,-5mm) {courant, amical};
  \node[anchor=north,font=\otRegular\fontsize{8}{10}\selectfont,text=otMuted] at (88mm,-5mm) {neutre, partout};
  \node[anchor=north,font=\otRegular\fontsize{8}{10}\selectfont,text=otMuted] at (130mm,-5mm) {travail, inconnus};
\end{tikzpicture}
\end{center}
\end{culture}

\begin{paire}
\begin{culture}[\en{How are you?} n'est pas une vraie question]
C'est une formule de salutation, comme « Ça va ? » quand on croise un collègue. On répond
court et positif — \en{Good, thanks. And you?} — même après une mauvaise nuit ! Vous
entendrez aussi \en{How are you doing?} et \en{How's it going?}
\end{culture}
\begin{culture}[La bise ? Jamais !]
Aux États-Unis comme au Royaume-Uni, on ne se fait pas la bise. Au travail, on se serre la
main (\en{a handshake}) lors de la première rencontre ; entre amis, un simple \en{Hi!}
suffit, parfois accompagné d'une accolade (\en{a hug}).
\end{culture}
\end{paire}

\begin{culture}[\en{Mr.}, \en{Mrs.} ou \en{Ms.} ?]
\begin{tabularx}{\linewidth}{@{}>{\enOn\bfseries}l<{\enOff}@{\hspace{3mm}}l@{\hspace{5mm}}Y@{}}
Mr.  & \ipa{ˈmɪstɚ} & « monsieur », pour tous les hommes : \en{Mr.\ Carter}\\
Mrs. & \ipa{ˈmɪsɪz} & « madame », pour une femme mariée : \en{Mrs.\ Lopez}\\
Ms.  & \ipa{mɪz}    & « madame », sans dire si elle est mariée : la forme la plus courante au travail\\
\end{tabularx}\par\vspace{1.6mm}
Toujours suivis du \textbf{nom de famille} : \en{Mrs.\ Lopez}, jamais \en{Mrs.\ Rosa}. Aux
États-Unis on écrit un point (\en{Mr.}) ; au Royaume-Uni, non (\en{Mr}). Pour s'adresser poliment à
un inconnu : \en{Excuse me, sir} (monsieur) ou \en{Excuse me, ma'am} (madame).
\end{culture}

%% ---------------------------------------------------------------------
\rubrique{pratique}{Exercices}

\exercice[ecrire,niveau=1]{Bonjour, à toute heure}{Regardez l'heure et la situation. Écrivez la
salutation qui convient : \en{good morning}, \en{good afternoon}, \en{good evening} ou \en{good night}.}
\begin{tcbraster}[raster columns=4,raster equal height=rows,raster column skip=3mm,raster before skip=1mm]
\otCarteHeure{8}{15}{8 h 15}{Léa arrive au studio.}{soleil}{Good morning!}
\otCarteHeure{14}{30}{14 h 30}{Léa entre dans un café.}{soleil}{Good afternoon!}
\otCarteHeure{19}{0}{19 h 00}{Léa arrive à un dîner chez Emma.}{lune}{Good evening!}
\otCarteHeure{23}{0}{23 h 00}{Léa quitte une fête pour aller dormir.}{lune}{Good night!}
\end{tcbraster}

\exercice[lire,niveau=1]{Question, réponse}{Associez chaque phrase (1 à 5) à la réponse qui convient (a à e).}
\begin{anglais}
\noindent\begin{tabularx}{\linewidth}{@{}>{\otSemi\fontsize{8.6}{10}\selectfont\color{otNavy}}l@{\hspace{2.4mm}}X@{\hspace{8mm}}>{\otSemi\fontsize{8.6}{10}\selectfont\color{otNavy}}l@{\hspace{2.4mm}}X@{}}
1. & Nice to meet you.          & a. & You're welcome.\\[1.4mm]
2. & How are you?               & b. & My name is Emma.\\[1.4mm]
3. & What's your name?          & c. & Bye! Have a nice evening!\\[1.4mm]
4. & Thank you very much!       & d. & Nice to meet you too.\\[1.4mm]
5. & See you tomorrow!          & e. & I'm fine, thanks. And you?\\
\end{tabularx}
\end{anglais}\par\vspace{3mm}
\noindent{\fontsize{10}{12}\selectfont\otSemi Réponses :}\hspace{4mm}%
1~→~\trou[9mm]{d}\hspace{6mm}2~→~\trou[9mm]{e}\hspace{6mm}3~→~\trou[9mm]{b}\hspace{6mm}%
4~→~\trou[9mm]{a}\hspace{6mm}5~→~\trou[9mm]{c}

\exercice[ecrire,niveau=1]{Le bon pronom}{Quel pronom remplace ces mots ? Écrivez \en{he},
\en{she}, \en{it}, \en{we}, \en{you} ou \en{they}.}
\begin{questions}[2]
\item Jack → \trou{he}
\item Emma → \trou{she}
\item Mr.\ Carter → \trou{he}
\item Mrs.\ Lopez → \trou{she}
\item the studio → \trou{it}
\item New York → \trou{it}
\item Jack and Emma → \trou{they}
\item Léa and I → \trou{we}
\item you and Jack → \trou{you}
\item Emma and Mrs.\ Lopez → \trou{they}
\end{questions}

\exercice[ecrire,niveau=1]{\en{am}, \en{is} ou \en{are} ?}{Complétez avec la bonne forme du verbe \en{be}.}
\begin{questions}[2]
\item I \trou[14mm]{am} Léa Martin.
\item You \trou[14mm]{are} Jack.
\item He \trou[14mm]{is} Mr.\ Carter.
\item My name \trou[14mm]{is} Emma.
\item She \trou[14mm]{is} very well.
\item We \trou[14mm]{are} fine, thanks.
\item Jack and Emma \trou[14mm]{are} in New York.
\item It \trou[14mm]{is} OK.
\end{questions}

\exercice[ecrire,niveau=2]{En version courte}{Récrivez chaque phrase avec la forme contractée.}
\begin{questions}[2]
\item I am Léa. → \trou[30mm]{I'm Léa.}
\item You are welcome. → \trou[30mm]{You're welcome.}
\item He is fine. → \trou[30mm]{He's fine.}
\item She is Emma. → \trou[30mm]{She's Emma.}
\item We are OK. → \trou[30mm]{We're OK.}
\item They are in New York. → \trou[30mm]{They're in New York.}
\item What is your name? → \trou[30mm]{What's your name?}
\item It is nice to meet you. → \trou[30mm]{It's nice to meet you.}
\end{questions}

\exercice[ecrire,niveau=2]{Dans le bon ordre}{Remettez les mots dans l'ordre pour former une
phrase. N'oubliez ni la majuscule, ni la ponctuation.}
\begin{questions}
\item \etiq{name} \etiq{Jack} \etiq{is} \etiq{my}\hfill\trou[66mm]{My name is Jack.}
\item \etiq{meet} \etiq{nice} \etiq{you} \etiq{to}\hfill\trou[66mm]{Nice to meet you.}
\item \etiq{you} \etiq{are} \etiq{how} \etiq{?}\hfill\trou[66mm]{How are you?}
\item \etiq{thanks} \etiq{I'm} \etiq{,} \etiq{fine}\hfill\trou[66mm]{I'm fine, thanks.}
\item \etiq{name} \etiq{your} \etiq{what's} \etiq{?}\hfill\trou[66mm]{What's your name?}
\item \etiq{you} \etiq{tomorrow} \etiq{see} \etiq{!}\hfill\trou[66mm]{See you tomorrow!}
\end{questions}

\exercice[ecrire,niveau=2]{\en{I'm} ou \en{my} ? \en{you're} ou \en{your} ?}{Complétez avec \en{I'm},
\en{my}, \en{you're} ou \en{your}. Attention aux majuscules en début de phrase.}
\begin{questions}[2]
\item Hello! \trou{My} name is Léa.
\item Hi! \trou{I'm} Jack.
\item What's \trou{your} name?
\item Thank you! — \trou{You're} welcome.
\item \trou{I'm} fine, thanks. And you?
\item Thank you, Emma. \trou{You're} very nice!
\end{questions}

\exercice[ecouter,niveau=1,piste=1.10]{Entendez-vous le \en{h} ?}{Écoutez la piste. Pour chaque
numéro, la voix dit un seul des deux mots, deux fois. Cochez le mot que vous entendez.}
\begin{questions}[4]
\item \choix{I}\quad\choix{hi}
\item \choix{and}\quad\choix{hand}
\item \choix{air}\quad\choix{hair}
\item \choix{eat}\quad\choix{heat}
\item \choix{old}\quad\choix{hold}
\item \choix{art}\quad\choix{heart}
\item \choix{ear}\quad\choix{here}
\item \choix{is}\quad\choix{his}
\end{questions}

\exercice[ecouter,niveau=2,piste=1.11]{Formel ou informel ?}{Écoutez les quatre mini-dialogues.
Pour chacun, cochez deux cases : la conversation est-elle formelle ou informelle ? Les personnes
se disent-elles bonjour ou au revoir ?}
\begin{center}\otTableau
\begin{tabular}{@{}l@{\hspace{10mm}}c@{\hspace{8mm}}c@{\hspace{14mm}}c@{\hspace{8mm}}c@{}}
 & \otTH{Formel} & \otTH{Informel} & \otTH{Bonjour} & \otTH{Au revoir}\\\midrule
{\otSemi Dialogue A} & \otCase & \otCase & \otCase & \otCase\\
{\otSemi Dialogue B} & \otCase & \otCase & \otCase & \otCase\\
{\otSemi Dialogue C} & \otCase & \otCase & \otCase & \otCase\\
{\otSemi Dialogue D} & \otCase & \otCase & \otCase & \otCase\\\bottomrule
\end{tabular}\end{center}

\exercice[ecouter,ecrire,niveau=2,piste=1.12]{Dictée à trous}{Le soir, Léa rencontre Sam, son
voisin. Écoutez et complétez le dialogue avec les mots qui manquent.}
\begin{anglais}
\noindent\begin{tabular}{@{}l@{\hspace{6mm}}l@{}}
\otNomDit{sam} & Hi! \trou{I'm} Sam.\\[1.8mm]
\otNomDit{lea} & Hello, Sam. My \trou{name} is Léa.\\[1.8mm]
\otNomDit{sam} & \trou{Nice} to meet you, Léa.\\[1.8mm]
\otNomDit{lea} & Nice to meet you \trou{too}.\\[1.8mm]
\otNomDit{sam} & How \trou{are} you?\\[1.8mm]
\otNomDit{lea} & I'm \trou{fine}, thanks. And you?\\[1.8mm]
\otNomDit{sam} & I'm OK. Have a nice \trou{day}!\\[1.8mm]
\otNomDit{lea} & Thank you. You \trou{too}!\\
\end{tabular}
\end{anglais}

\exercice[ecrire,niveau=2]{Traduisez}{Traduisez en anglais. Plusieurs réponses sont parfois possibles.}
\begin{questions}[1][french]
\item Bonjour, je m'appelle Paul. \textit{(le matin)}\lignes{1}
\item Enchanté !\lignes{1}
\item Comment allez-vous ? — Très bien, merci. Et vous ?\lignes{1}
\item Salut, Emma ! À demain ! \textit{(en partant)}\lignes{1}
\item Bonsoir, madame Lopez. \textit{(en arrivant)}\lignes{1}
\item Merci beaucoup. — De rien.\lignes{1}
\item Je te présente Jack.\lignes{1}
\item Ils sont à New York.\lignes{1}
\end{questions}

\exercice[ecouter,parler,niveau=3,piste=1.13]{À vous de parler !}{Jack vous parle. Répondez-lui
\textbf{à voix haute} pendant chaque pause, en vous aidant des indications. Comparez ensuite
avec les réponses modèles du corrigé, puis recommencez sans regarder le livre.}
\begin{center}\otTableau
\begin{tabular}{@{}>{\otSemi\fontsize{8.6}{10}\selectfont\color{otNavy}}l@{\hspace{3mm}}>{\enOn}p{70mm}<{\enOff}@{\hspace{6mm}}p{78mm}@{}}
 & \otTH{Jack dit\ldots} & \otTH{Vous\ldots}\\\midrule
1. & Hi! I'm Jack. What's your name? & dites bonjour et donnez votre prénom.\\
2. & Nice to meet you! & répondez poliment.\\
3. & How are you? & dites que vous allez bien et demandez-lui comment il va.\\
4. & I'm good, thanks. Welcome to New York! & remerciez-le.\\
5. & OK, see you tomorrow! & dites au revoir et souhaitez-lui une bonne soirée.\\\bottomrule
\end{tabular}\end{center}

%% ---------------------------------------------------------------------
\rubrique{bilan}{Faites le point}

\begin{essentiel}
\begin{minipage}[t]{0.56\linewidth}
\otEssTitre{Les phrases clés}
\renewcommand{\arraystretch}{1.32}
\begin{tabular}{@{}>{\enOn\otSemi\fontsize{9.4}{11}\selectfont}l<{\enOff}@{\hspace{3mm}}>{\fontsize{9.2}{11}\selectfont\color{white!78!otNavy}}l@{}}
Hello! / Hi!                 & Bonjour ! / Salut !\\
Good morning.                & Bonjour (le matin).\\
Good evening.                & Bonsoir (en arrivant).\\
I'm Léa. / My name is Léa.   & Je m'appelle Léa.\\
What's your name?            & Comment t'appelles-tu ?\\
Nice to meet you.            & Enchanté(e).\\
How are you?                 & Comment ça va ?\\
I'm fine, thanks. And you?   & Bien, merci. Et toi ?\\
Goodbye! / See you tomorrow! & Au revoir ! / À demain !\\
Thank you. — You're welcome. & Merci. — De rien.\\
\end{tabular}
\end{minipage}\hfill
\begin{minipage}[t]{0.38\linewidth}
\otEssTitre{Le verbe be}
\renewcommand{\arraystretch}{1.32}
\begin{tabular}{@{}>{\enOn\fontsize{9.4}{11}\selectfont}l<{\enOff}@{\hspace{2mm}}>{\enOn\otSemi\fontsize{9.4}{11}\selectfont}l<{\enOff}@{}}
I am      & → I'm\\
he is     & → he's \textmd{(she's, it's)}\\
you are   & → you're \textmd{(we're, they're)}\\
\end{tabular}
\otEssTitre{Les pronoms sujets}
{\enb{\otSemi\fontsize{9.6}{12}\selectfont I · you · he · she · it · we · you · they}}\par
{\fontsize{9}{11}\selectfont\color{white!78!otNavy}\en{I} toujours en majuscule ; \en{you} = tu et vous}
\otEssTitre{Prononciation}
{\fontsize{9.2}{11.6}\selectfont \textbf{h} = un souffle : \en{\textbf{h}ello, \textbf{h}i, \textbf{h}ow}.\par
Pas de h devant \en{I'm}, \en{and}, \en{eat} !\par}
\end{minipage}
\end{essentiel}


\begin{jesais}
\item saluer quelqu'un et prendre congé, à chaque moment de la journée ;
\item me présenter (\en{I'm…}, \en{My name is…}) et demander son nom à quelqu'un ;
\item réagir à une présentation : \en{Nice to meet you (too).} ;
\item demander comment ça va, et répondre brièvement ;
\item utiliser \en{I, you, he, she, it, we, they} et conjuguer \en{be} au présent ;
\item prononcer le \textbf{h} de \en{hello}, \en{hi}, \en{how}… sans en ajouter devant \en{I'm} ou \en{and}.
\end{jesais}

\clearpage
\begin{testexpress}
\qcm[a]{Vous arrivez au bureau à 9 heures. Vous dites :}{\en{Good morning.}}{\en{Good night.}}{\en{Goodbye.}}
\qcm[c]{« Enchanté ! » se dit :}{\en{Welcome!}}{\en{You're welcome!}}{\en{Nice to meet you!}}
\qcm[b]{« Ils sont » se dit :}{\en{they is}}{\en{they are}}{\en{them are}}
\qcm[b]{Complétez : \en{She \rule[-0.6mm]{9mm}{0.5pt} Emma.}}{\en{am}}{\en{is}}{\en{are}}
\qcm[a]{La forme contractée de \en{you are} est :}{\en{you're}}{\en{your}}{\en{yo're}}
\qcm[c]{« Comment tu t'appelles ? » se dit :}{\en{How are you?}}{\en{Who are you?}}{\en{What's your name?}}
\qcm[b]{Le pronom \en{it} désigne :}{une femme}{une chose ou un animal}{un homme}
\qcm[c]{Dans quel mot le \en{h} ne se prononce-t-il pas ?}{\en{hello}}{\en{how}}{\en{hour}}
\qcm[a]{La réponse naturelle à \en{How are you?} est :}{\en{Fine, thanks. And you?}}{\en{I'm Léa.}}{\en{Nice to meet you.}}
\qcm[c]{On dit \en{Good night} :}{en arrivant le soir}{le matin}{en partant le soir, ou avant de dormir}
\end{testexpress}

\noindent{\fontsize{9.6}{12}\selectfont\color{otMuted}Niveau visé (CECRL, niveau A1) : « peut utiliser
des formules simples de salutation et de prise de congé, se présenter et répondre à des questions
simples sur soi ». Tous les corrigés de l'unité sont page~\pageref{cor-u1-1} ; les textes des
pistes d'exercices, page~\pageref{trans-1.10}.\par}
